%%%%%%%%%%%%%%%%%%%%%%%%%%%%%%%%%%%%%%%%%%%%%%%%%%%%%%%%%%%%%%%%%%%%%%%%%%%%%%%%
%%  content/06-conclusion.tex — فصل ۶: نتیجه‌گیری و کارهای آینده
%%%%%%%%%%%%%%%%%%%%%%%%%%%%%%%%%%%%%%%%%%%%%%%%%%%%%%%%%%%%%%%%%%%%%%%%%%%%%%%%
\section{نتیجه‌گیری و کارهای آینده}

%%%%%%%%%%%%%%%%%%%%%%%%%%%%%%%%%%%%%%%%%%%%%%%%%%%%%%%%%%%%%%%%%%%%%%%%%%%%%%%%
\begin{frame}{جمع‌بندی}
  \begin{itemize}
    \item سامانه‌ای ساخته شد که پرسش زبان طبیعی را می‌گیرد و پاسخی
          \hl{مستند} با ارجاع به منبع بازمی‌گرداند.
    \item زنجیرهٔ \term{بازیابی، بازرتبه‌بندی و تولید} به اجزای جایگزین‌پذیر شکسته
          شد، پس هر مرحله جداگانه سنجیدنی است.
    \item \term{بازرتبه‌بندی دومرحله‌ای} بافتار ورودی مدل را به سه سند کوتاه کرد و
          در همان حال دقت بازیابی را بالا برد.
  \end{itemize}

  \begin{PTresult}
    درستی پاسخ \hl{۹۵٪} در برابر \hl{۷۷٫۵٪} سامانهٔ مرجع، با بافتاری کوتاه‌تر و
    هزینهٔ کمتر برای هر پرسش.
  \end{PTresult}
\end{frame}

%%%%%%%%%%%%%%%%%%%%%%%%%%%%%%%%%%%%%%%%%%%%%%%%%%%%%%%%%%%%%%%%%%%%%%%%%%%%%%%%
\begin{frame}{محدودیت‌ها}
  \begin{PTcons}
    \begin{itemize}
      \item وابستگی به سرویس‌های بیرونی: هزینه، سهمیهٔ فراخوانی و دسترس‌پذیری
            خارج از کنترل سامانه است.
      \item سه سند نهایی برای پرسش‌هایی که پاسخشان در وب پراکنده است کم است.
      \item مجموعهٔ آزمون \hl{۴۰ پرسشی} کوچک است و نتیجه را باید مقدماتی دانست.
    \end{itemize}
  \end{PTcons}

  \begin{itemize}
    \item هیچ‌یک از این محدودیت‌ها به \term{معماری} گره نخورده است؛ هر سه با
          تعویض یا تنظیم یک جزء برطرف می‌شوند.
  \end{itemize}
\end{frame}

%%%%%%%%%%%%%%%%%%%%%%%%%%%%%%%%%%%%%%%%%%%%%%%%%%%%%%%%%%%%%%%%%%%%%%%%%%%%%%%%
\begin{frame}{کارهای آینده}
  \begin{enumerate}
    \item \term{تعداد سند پویا}: انتخاب \(k\) بر پایهٔ امتیاز بازرتبه‌بند، به‌جای
          عدد ثابت سه.
    \item \term{بازرتبه‌بند بومی}: ریزتنظیم یک مدل رمزگذار متقابل روی دادهٔ فارسی
          برای حذف وابستگی به سرویس بیرونی.
    \item \term{حافظهٔ گفت‌وگو}: پشتیبانی از پرسش‌های پیاپی و ارجاع‌های درون‌متنی.
    \item \term{ارزیابی گسترده‌تر}: مجموعهٔ چند‌صد پرسشی با داوری چنددادوری و
          گزارش بازهٔ اطمینان.
  \end{enumerate}

  \begin{PTnote}
    ترتیب بالا بر پایهٔ نسبت \hl{اثر به هزینه} چیده شده است؛ گام نخست کمترین
    تغییر در کد و بیشترین اثر را دارد.
  \end{PTnote}
\end{frame}
